\begin{toreview}

\begin{abstract}
Alignment faking occurs when a model strategically complies with a training objective to avoid changes to its existing behavior. Whether evidence obtained from individual models generalizes across smaller language models remains unclear. Motivated by a systematic review of 80 entries in AI safety research, this question is explored through three prompt-based experiments: harmful requests, persona preservation, and animal welfare dismissal. The experiments compare simulated free-tier interactions described as eligible for retraining with paid-tier interactions described as excluded from retraining, with and without explicit scratchpad prompting. We evaluate models from five families using a local response classifier, self-reported acceptance tags, and manually reviewed scratchpad classifications. In the harmful-request replication, LLaMA 3.1 8B without scratchpad prompting produces 11 classifier-judged compliant responses out of 100 in the free tier and none in the paid tier. The two-sided Fisher's exact test gives $p=0.00073$, which remains significant under a Bonferroni correction across 28 harmful-request comparisons. With scratchpad prompting, both rates are 5\%. No other harmful-request comparison reaches the nominal 0.05 threshold. Persona preservation and animal welfare dismissal provide no consistent behavioral evidence extending this result, while self-reported acceptance frequently exceeds classifier-judged compliance. These findings reproduce a specific compliance gap but do not establish alignment faking as a general property of the tested models. They also show why behavioral differences, acceptance tags, and tier-aware reasoning should be treated as distinct evaluation signals.
\end{abstract}

\end{toreview}